\begingroup\small\setlength{\tabcolsep}{4pt}
\begin{longtable}{@{}p{0.04\linewidth}>{\raggedright\arraybackslash}p{0.34\linewidth}>{\raggedright\arraybackslash}p{0.56\linewidth}@{}}
\toprule
No. & Delivery & Comparison\\
\midrule\endhead
\bottomrule\endfoot
1 & \file{pa_unary_power_cover} & \textbf{Unary theory}. Unary power-covering theory calibrated to PA by explicit finite-proof transformations; includes a zero-free separation model and a dyadic implementation. No absolute smallest wrapper or new physical TM saving is claimed.\\
\addlinespace[6pt]
2 & \file{U208_ConPA_Conditional_Context_2026-10-29} & \textbf{U208 contextual quotient}. Eliminates native 172 using the predecessor contexts of 172 and 211. The source is unchanged. Main/compatibility parents yield the two U208 graph classes; alternate helper choice is included.\\
\addlinespace[6pt]
3 & \file{U208_ConPA_Entry_Context_2026-10-29} & \textbf{U208 contextual quotient}. The same 172/211 saving is presented as entry-context elimination; an all-initial-tape local simulation and unchanged source edges are supplied. This is not a second additive saving.\\
\addlinespace[6pt]
4 & \file{U208_ConPA_History_2026-10-29} & \textbf{U208 contextual quotient}. The same 172/211 saving is presented with a finite-history proof; four arbitrary-exterior macro cases establish the local simulation. A source-preserving rotation is separately rejected by its physical contracts.\\
\addlinespace[6pt]
5 & \file{U208_ConPA_2026-10-29} & \textbf{U208 contextual quotient}. Three-cell presentation of the same contextual fold, deleting 211 and redirecting it to 172. Dense relabeling resolves the native naming difference; ordinary mathematical proof, not a new end-to-end kernel theorem.\\
\addlinespace[6pt]
6 & \file{u208_research_2026-10-29} & \textbf{U208 contextual quotient}. Two 208-state constructions, corresponding to the two supplied U209 classes, with complete symbolic runtime/startup replays. They implement the same contextual saving.\\
\addlinespace[6pt]
7 & \file{Below_U208_Research_2026-10-29} & \textbf{U208 negative study}. New full-frame boundary projections and closed-loop source pruning strengthen local searches. No 207-state machine is established; nine principal 208 presentations still form two classes.\\
\addlinespace[6pt]
8 & \file{u208_comparison_research_2026-10-29} & \textbf{U208 negative study}. Rechecks both 208 classes, source/macro/table search templates, and an explicit failed 207 repair. A source DAG size or finite screen is not a physical reduction certificate.\\
\addlinespace[6pt]
9 & \file{ConPA_below208_research_2026-10-29} & \textbf{U208 negative study}. Shows a role-identification weakening of the proposed unary theory, and a restricted 197-node source-DAG minimum for a fixed family. Synthesizes both keeper rows over all state pairs in stated bounded windows; no 207 construction.\\
\addlinespace[6pt]
10 & \file{pa_consistency_research_2026-10-29} & \textbf{U208 source specialization}. Reduces the source to 3,904 primitives using inductive register abstractions from the canonical seed. The source theorem is not compiled to an accepted smaller binary table; rejects a 207 rotation experiment.\\
\addlinespace[6pt]
11 & \file{pa208_research_2026-10-29} & \textbf{U208 helper variants}. Four local merge/helper constructions supply 18 additional 208-state graphs, nine over each original class. Their exact incoming contexts and local simulations are given; no 207 result.\\
\addlinespace[6pt]
12 & \file{pa_tm_research_2026-11-13} & \textbf{PA196 layout}. Reorders axiom heads as(1,2,3,5,6,4), reducing maximum source path 38\ensuremath{\to}36 while retaining 196 states. Separately rebuilds the published physical Lean coverage; the new layout has symbolic/startup checks, not a new kernel calibration.\\
\addlinespace[6pt]
13 & \file{pa_consistency_research_2026-11-13} & \textbf{Diagonal-vacuity calculus}. Introduces a shorter 191-token diagonal-vacuity emitter: 3,700 primitive sites, verified 199-state table. Its paper-level PA calibration and logical witness interfaces are distinct from the state-count result.\\
\addlinespace[6pt]
14 & \file{conpa_tm_comparison_and_reduction_search} & \textbf{PA196 command experiments}. Two one-row MP controller changes with mathematical command arguments and bounded frame tests. Their own packet labels them experimental; two bad Shift layouts halt after 43,563 blank transitions.\\
\addlinespace[6pt]
15 & \file{pa196_research_2026-11-13} & \textbf{Source factorization}. A157-internal-node representation preserves all 3,812 primitive instructions and source edges. The prospective 195-state price loses a necessary reader helper and is invalid; a source DAG node is not a TM state.\\
\addlinespace[6pt]
16 & \file{pa_consistency_research} & \textbf{Source factorization and scanner}. A158-node source factorization is compiled to 196 states: an extra reader fold is offset by an eight-zero scanner. Supplies rewrite-head countermodels and logical validation of 5,760 emitter variants, not 5,760 physical machines.\\
\addlinespace[6pt]
17 & \file{pa_tm_comparison_and_reductions} & \textbf{PA196 layouts and logical rewrite}. Two new 196 layouts pass full runtime/startup checks. A199-token emitter with 3,773 primitives compiles to 198 states; smaller grammar witnesses alone do not prove a reduction.\\
\addlinespace[6pt]
18 & \file{pa_tm_extended_research_2026-11-13} & \textbf{PA196 invariants and closing rows}. Four one-row 196 variants plus joint source invariants reducing reachable outcomes 4,959\ensuremath{\to}4,954. Gives a countervaluation for a shorter quantifier-elimination head and records the missing atomic-vacuity obligation.\\
\addlinespace[6pt]
19 & \file{conpa_tm_comparison_and_search} & \textbf{PA196 physical placement}. Swaps the physical A/D registers to obtain a new 196 table with clean startup. A source-preserving 195 rotation fails 93 primitive contracts; a reflexivity-only equality shortcut loses the required logical strength.\\
\addlinespace[6pt]
20 & \file{pa_consistency_reduction} & \textbf{PA196 and integrated witness gate}. Rebracketed and cached-only-Shift 196 variants are fully bound. An integrated equality-witness gate shortens the source to 3,745 sites but its best checked compilation here uses 203 controls.\\
\addlinespace[6pt]
21 & \file{pa_consistency_comparison_2026_11_13} & \textbf{Split quantifier calculus}. Splits the first quantifier head into a modal-B form plus explosion, with finite Hilbert proofs in both directions. Contains composition/grammar experiments but no accepted smaller binary machine.\\
\addlinespace[6pt]
22 & \file{PA_consistency_TM_research_2026-11-13} & \textbf{ZFC-inspired source experiments}. Provides a 157-node representation of a 3,573-site combined source, an original-source zero-run-five grammar, and a parameter-free residual-MP command design adapted from ZFC. No accepted 195 table results.\\
\addlinespace[6pt]
23 & \file{pa_consistency_research_2026-11-13_2} & \textbf{PA196/198 layouts and MP halt}. A rotated 196 layout and a five-fold 198 implementation of diagonal vacuity pass checks. A command-level variant halts immediately when MP derives 1; this packet does not bind that idea to a new binary TM.\\
\addlinespace[6pt]
24 & \file{conpa_tm_study_2026-11-13} & \textbf{Cached-controller composition}. Combined 196 joins cached-result MP with cached-only Shift, with 4,928 arbitrary-register outcomes and a written calibration. Also includes direct-negation 197 and a 207 physical realization of a repaired 157-node source.\\
\addlinespace[6pt]
25 & \file{pa_below196_research} & \textbf{Logical and physical negative study}. Studies split gates, transposed implication pairing, source grammars and helper substitutions. Reversing implication arguments has a 37-site macro but raises the shared-DAG cost; no complete transposed logical/physical construction is supplied.\\
\addlinespace[6pt]
26 & \file{pa_consistency_comparison_and_new196_2026-11-13} & \textbf{Separate-diagonal family}. Three 196 layouts include the 3,700-site diagonal emitter and the 3,673-site separate-diagonal emitter. The latter 190-token source is the ancestor of the later 195/194 frontier.\\
\addlinespace[6pt]
27 & \file{pa_tm_research_2026-11-13_2} & \textbf{Known-zero jump and parity}. A source-proved zero decrement replaces a jump at site 26, producing a distinct 196 table. A 195 fold fails the source binding; a 202-state base-eight parity runtime lacks a certified initializer.\\
\addlinespace[6pt]
28 & \file{pa_consistency_cross_archive} & \textbf{Cross-archive controller combinations}. Four 196 tables combine diagonal/separate emitters with the cached controller and optional known-zero jump; all have full runtime/startup bindings. A separate shorter-negation composition remains 197.\\
\addlinespace[6pt]
29 & \file{pa195_conpa_2026-11-13} & \textbf{PA195 association}. Uses the separate-diagonal source unchanged with 68 fresh readers and seven compatible folds: 134+68\ensuremath{-}7=195. Includes the complete local quotient proof and source/physical/startup bindings.\\
\addlinespace[6pt]
30 & \file{pa_consistency_reduction_audit_2026-11-14} & \textbf{PA196 composition audit}. Four compositions of diagonal/separate emitters and cached-result MP yield 196/197 tables. Short-negation compositions still require a global scratch-zero invariant for their claimed command calibration.\\
\addlinespace[6pt]
31 & \file{pa_consistency_research_2026-11-14} & \textbf{NEG18}. Replaces eight 47-site encoded-negation bodies by 18-site implementations. The 3,441-site source and zero-scratch invariant are checked, but physical cost is 134+68\ensuremath{-}5=197.\\
\addlinespace[6pt]
32 & \file{PA195_ConPA_research_2026-11-14} & \textbf{PA195 association}. Primary 195S uses 67 readers and six folds with the separate source. Supplemental 195C uses the combined controller and 3,649 sites; this is a different allocation from the 68/7 PA195.\\
\addlinespace[6pt]
33 & \file{pa195_followup_research} & \textbf{PA195 scanner tradeoff}. M195S8 uses an eight-zero scanner and eight folds from 203 controls. The tempting 194 short-scanner quotient passes local folds/CFG but fails a legal seven-zero source address.\\
\addlinespace[6pt]
34 & \file{pa194_conpa_2026-11-13} & \textbf{Primary PA194}. Separate-diagonal source with 67 readers and seven folds: 134+67\ensuremath{-}7=194. Supplies original and A/D physical-register placements; the source calculus and original sum Shift are unchanged.\\
\addlinespace[6pt]
35 & \file{pa194_conpa_2026-11-13_2} & \textbf{Alternate PA194}. Another association of the same separate-diagonal source uses 68 readers and eight folds: 134+68\ensuremath{-}8=194. The literal graph differs from the 67/7 construction.\\
\addlinespace[6pt]
36 & \file{pa195_diagonal_combined_2026-11-13} & \textbf{PA195 diagonal-controller composition}. PA195DC combines the diagonal-vacuity emitter and cached-controller changes at 195 states. It is a larger but semantically different antecedent of cached 194 designs.\\
\addlinespace[6pt]
37 & \file{pa194C_conpa_2026-11-13_2} & \textbf{Cached PA194C}. One of five deliveries of the same PA194C graph: MP caches its result and Shift tests only the cached head. Requires an adapted command compiler; there is no new sub 194 count.\\
\addlinespace[6pt]
38 & \file{pa194_new_designs_2026-11-13} & \textbf{Cached PA194C/CS}. PA194C duplicates the cached graph; PA194CS additionally uses a symmetric destructive comparison to remove a transfer loop. Both remain 194 with separately bound controllers.\\
\addlinespace[6pt]
39 & \file{pa194c_conpa} & \textbf{Cached PA194C}. Another delivery of the same cached-head PA194C literal graph and its mathematical proof. File/report multiplicity is not construction multiplicity.\\
\addlinespace[6pt]
40 & \file{pa194c_conpa_2026-11-13_1} & \textbf{Cached PA194C}. Another delivery of the same cached-result-MP/cached-only-Shift PA194C graph and proof bindings.\\
\addlinespace[6pt]
41 & \file{pa194c_conpa_2026-11-13} & \textbf{Cached PA194C}. Another delivery of the same PA194C graph, with inherited calibration and executable source/physical evidence.\\
\addlinespace[6pt]
42 & \file{pa194_radix_family_2026-11-13} & \textbf{Coefficient family}. M5/M6 replace the formula-constructor coefficient 4 by 5/6 with uniform fresh-label PA calibration. Also supplies M4R, a same-coefficient reassociation; coefficient changes require fresh source/table bindings.\\
\addlinespace[6pt]
43 & \file{pa194j_pa194d_conpa_2026-11-13} & \textbf{Known-zero and diagonal head}. PA194J replaces ten jumps with source-proved zero decrements. PA194D changes the separate diagonal formula from \ensuremath{\Gamma}(A) to\ensuremath{\Gamma}(D), with revised selectors and a 69-reader/nine-fold realization.\\
\addlinespace[6pt]
44 & \file{pa194q_conpa_2026-11-13} & \textbf{Quantifier exchange and zero return}. PA194Q exchanges curried premises of the universal-distribution head, justified by finite K/S/MP exchange. PA194CSZ equals the previously named PA194X known-zero cached-symmetric graph.\\
\addlinespace[6pt]
45 & \file{PA194R_research_2026-11-14} & \textbf{Reassociated zero-closing decoder}. PA194F uses a zero-proved final AX decrement; PA194R reassociates that source to a 67-reader decoder, changing 1,111 source addresses. Arithmetic semantics are retained with new physical bindings.\\
\addlinespace[6pt]
46 & \file{pa194N_conpa_2026-11-13} & \textbf{No-push MP}. PA194N omits the dummy-zero push on successful MP and leaves its result in the visible top, with the original sum Shift. Completeness uses a reserved-zero proof compiler valid for every incoming store.\\
\addlinespace[6pt]
47 & \file{pa194zr_conpa_2026-11-13} & \textbf{Zero-return convention}. PA194ZR replaces eight comparator returns and the AX exit by zero-decrement returns, under a 129-site invariant. The initializer is checked independently; apparently similar patches can fail startup.\\
\addlinespace[6pt]
48 & \file{pa194hs_conpa_2026-11-13} & \textbf{No-push cached Shift}. PA194H combines visible-result no-push MP with cached-only Shift; PA194HS makes that cached-only comparison symmetric. HS is the graph also called NS by the older lowercase packet.\\
\addlinespace[6pt]
49 & \file{pa194n5_n6_conpa_2026-11-13} & \textbf{No-push coefficient family}. PA194N5/N6 combine coefficient 5/6 formula constructors with no-push MP and the original sum Shift. N6 is delivered again in a later packet with the same table.\\
\addlinespace[6pt]
50 & \file{pa194ns_conpa_2026-11-13} & \textbf{No-push symmetric cached Shift}. This older lowercase PA194NS is byte-identical to PA194HS: cached-only symmetric Shift. It must be distinguished from the later uppercase PA194NS\_research sum-Shift machine.\\
\addlinespace[6pt]
51 & \file{pa194n6_research_2026-11-13} & \textbf{No-push coefficient family}. A repeat delivery of PA194N6, with coefficient 6 constructors, no-push MP and original sum Shift; the same graph already appears in the N5/N6 family.\\
\addlinespace[6pt]
52 & \file{PA194NS_research} & \textbf{Symmetric original sum Shift}. PA194SUMS symmetrizes the original sum comparison while retaining dummy-push MP; the new PA194NS adds no-push MP. This NS is different from the older cache-only NS=HS.\\
\addlinespace[6pt]
53 & \file{pa194haz_research_2026-11-14} & \textbf{Cached result with register swap}. PA194HA combines cached-result MP, original sum Shift and A/D physical placement; HAZ adds a final source-proved zero decrement. The packet explicitly specifies native start 74 and dense start 0; these label systems must not be mixed.\\
\addlinespace[6pt]
54 & \file{pa194qn_research_2026-11-13} & \textbf{Q plus no-push}. PA194QN combines the exchanged distribution head with no-push MP and the original sum Shift, at coefficient 4. The later PA194NQ delivery is byte-identical.\\
\addlinespace[6pt]
55 & \file{pa194h6_conpa_2026-11-13} & \textbf{No-push cached coefficient6}. PA194H6 combines coefficient 6 formula arithmetic, no-push MP and cached-only Shift. It changes a command halt condition relative to N6, not merely a physical state naming.\\
\addlinespace[6pt]
56 & \file{pa194dhs_conpa_2026-11-13} & \textbf{Diagonal plus cached symmetry}. PA194DHS combines \ensuremath{\Gamma}(D), no-push MP and symmetric cached-only Shift; its count is 135+69\ensuremath{-}10. Its selectors/compiler differ from the \ensuremath{\Gamma}(A) family.\\
\addlinespace[6pt]
57 & \file{pa194nq_research_2026-11-13} & \textbf{Q plus no-push; rejected radix3}. PA194NQ repeats PA194QN. A 193-state factor 3 candidate passes physical checks but is logically rejected because constructor collision yields a derivation of contradiction.\\
\addlinespace[6pt]
58 & \file{pa194_qns_research_2026-11-13} & \textbf{Q plus symmetric sum Shift}. PA194QNS combines Q, no-push MP and symmetric original sum Shift; PA194QS retains dummy-push MP. Also proves a U/E constructor-tag interchange, without supplying its literal TM.\\
\addlinespace[6pt]
59 & \file{pa194_qradix_research_2026-11-13} & \textbf{Q coefficient family}. PA194Q5/Q6 combine Q with coefficient 5/6, retaining original dummy-push MP and original sum Shift. These are the direct parents of the new QN5/QN6 composition.\\
\addlinespace[6pt]
60 & \file{pa194qh6_conpa_2026-11-13} & \textbf{Q no-push cached coefficient6}. PA194QH6 combines Q, coefficient 6, no-push MP and cached-only Shift. New QN6 instead retains the original sum Shift; their halting guards are different.\\
\addlinespace[6pt]
61 & \file{pa194qns_research_2026-11-13} & \textbf{Q plus symmetric sum Shift}. Another delivery of the identical PA194QNS graph, using no-push MP and symmetric original sum Shift. This delivery does not add another graph class.\\
\addlinespace[6pt]
62 & \file{pa194qh_conpa_2026-11-13.zip} & \textbf{Q no-push cached Shift}. New PA194QH: exchanged distribution head, no-push MP and cached-only Shift at coefficient 4. Reports 194 states, 3,653 source sites, 4,754 arbitrary-register outcomes, and 57,922-step startup. Same architecture as QH6 but original coefficient 4.\\
\addlinespace[6pt]
63 & \file{pa_consistency_reduction_study_2026-11-13.zip} & \textbf{Restricted decoder lower bound}. Negative study: an exact rational dual proves no saving beyond seven controls in a fixed raw 201-parent/fresh-direct-selector/disjoint-reader-fold template. Source reassociation removes one fixed scanner obstruction, but provides no smaller initialized TM. Includes an unvalidated N6 append-zero mutation.\\
\addlinespace[6pt]
64 & \file{pa_consistency_research_2026-11-14.zip} & \textbf{All-offset radix3 obstruction}. Negative study proves that every nonnegative constant equality offset fails in the specified radix 3 constructor scheme with retained heads/rules. Records mixed 3/6 and split-gate failures, plus 23 literal source-rotation inputs at 194 states.\\
\addlinespace[6pt]
65 & \file{pa_sub194_research.zip} & \textbf{MP-success halt at196}. MP196 halts immediately when successful MP derives 1; other success leaves the consequent visible without pushing. A source invariant excludes the old Shift halt. Reports 196 states, 3,663 sites, 4,768 outcomes, 58,154-step startup.\\
\addlinespace[6pt]
66 & \file{pa_sub194_research_2026-11-13.zip} & \textbf{193 runtime without startup}. A 193-control runtime folds 301 into 88 with two row changes and passes all 4,780 arbitrary-register outcomes. Its blank start enters an outward ray; 1,158 elementary startup variants are rejected. Also reports checked 194-state F\_AD/R\_AD companions and a broader admissible-seed theorem.\\
\addlinespace[6pt]
\end{longtable}
\endgroup
